% title:          Zero field points of charged triangle
% area:           electromagnetism
% methods:
% difficulty:
% difficulty-ai:  medium
% status:
% review:         none
% --- history: all optional, leave EMPTY if unknown ---
% origin:         folklore
% origin-date:
% origin-number:
% also-in:
% taken-from:
% references:     Classical special case of Maxwell's problem on equilibrium points of point charges: Maxwell, A Treatise on Electricity and Magnetism, section 113 (at most (l-1)^2, i.e. 4 for three charges), as quoted in A. Gabrielov, D. Novikov, B. Shapiro, Mystery of point charges, arXiv 2004 - https://arxiv.org/pdf/math-ph/0409009; K. Killian, A remark on Maxwell's conjecture for planar charges, Complex Var. Elliptic Eqns 54 (2009) 1073-1078: exactly four equilibrium points for unit charges at the vertices of an equilateral triangle (as cited by Bilogliadov); M. Bilogliadov, Equilibria of the field generated by point charges, arXiv 2014, Examples 3.2 and 5.4 (the centre plus one point on each perpendicular bisector; for the Coulomb case, a root of r^5+5r^4+r^3+r^2-4r+1 with circumradius 1) - https://arxiv.org/pdf/1404.7198; Physics Stack Exchange Q108929 (2014), Electric Field inside a regular polygon with corner charges - https://physics.stackexchange.com/questions/108929
% history:        ai
\begin{problem}
There is an equilateral triangle with equal charges kept on its vertices. Find the number and location of points with electric field $\vec{E}_{total}=0$.
\\ For a point charge q,
 $$\vec{E}(\vec{r})=\frac{q(\vec{r}-\vec{r}_0)}{4\pi\epsilon_0 |\vec{r}-\vec{r}_0|^3} $$
where $\vec{r}_0$ is the location of the point charge. Take $q=1$ and $4\pi\epsilon_0=1$. (The side length of the triangle can be chosen arbitrarily).
\end{problem}
